\section{Measuring Self-Recognition}
\label{sec:measure}

A \emph{panel} is a set of models that both write texts and judge them, so
each judge sees some texts of its own among the texts it judges. For each
judge we compute three rates from the same set of judgments:
\begin{itemize}
\item \textbf{self-naming rate}: the share of the judge's own texts that it
says it wrote;
\item \textbf{false-alarm rate}: the share of other models' texts that it
says it wrote;
\item \textbf{peer baseline}: the share of the judge's texts that the
\emph{other} judges attribute to it.
\end{itemize}
From these we take two differences. \textbf{Discrimination} is the
self-naming rate minus the false-alarm rate: does the judge claim its own
text more than others'? \textbf{Self-advantage} is the self-naming rate minus
the peer baseline: does the judge know its own text better than other judges
do? With five candidate authors, guessing at random gives 20\% for all three
rates and zero for both differences.

\paragraph{Two tests} The \emph{usual test} asks only whether the
self-naming rate is above chance. \emph{Our test} requires both a
self-advantage and discrimination. Discrimination alone already means
anonymisation fails for that judge; the self-advantage shows whether the
judge has any special knowledge of its own writing, beyond what any judge
could see.

\paragraph{An example} In the lineup (Section~\ref{sec:frontier}), Claude
names itself on 33 of its 38 texts (87\%) and on only 14 of the 152 texts it
did not write (9\%). The other four judges name Claude on 54\% of Claude's
texts. Claude discriminates and beats its peers by 33 points, so it passes.
Shown the same texts one at a time, Claude still names itself on 33 of 38.
But now its peers name Claude on 82\% of Claude's texts, because in this
format every judge names Claude or GPT for almost everything. Claude's
self-naming rate is unchanged, yet its self-advantage nearly disappears. In
the same format GPT names itself on 36 of its 38 texts, the highest rate in
the study, but also on 79 of the 152 texts it did not write. The self-naming
rate alone cannot tell these cases apart.

\paragraph{Statistics} Judgments of the same prompt are not independent,
since every judge sees the same few texts. All our tests therefore resample
or shuffle whole prompts or whole texts, never single judgments: 95\%
intervals come from resampling prompts, and $p$-values from permutation
tests (shuffling which judge counts as ``self'' on each text, or which text
in a lineup is the judge's own). We correct for testing several judges at
once with Holm's method. For yes/no questions answered with a probability,
we report the AUC, which does not depend on how readily a judge says yes.
Appendix~\ref{app:stats} gives details.
